\subsection{1. Task 1 --- AA for Gibbs and Dobrushin measures}\label{task-1-aa-for-gibbs-and-dobrushin-measures}

\subsubsection{1A. Formal targets}\label{a.-formal-targets}

For a Boolean Gibbs measure \(\mu\) with bit-level Dobrushin coefficient \(\alpha\le1-\delta\), bounded range/degree/fields, and every multilinear polynomial \(p:\{0,1\}^n\to[0,1]\) of degree at most \(d\), the AA goal is

\[
\max_i\operatorname{Inf}_i^\mu(p)
\ge c(\delta,\delta_0,\Delta,d,R_0)
\left(\frac{\operatorname{Var}_\mu(p)}d\right)^{C(\delta,\Delta,R_0,d)},
\]

where \(c,C>0\) are independent of \(n\). Here

\[\operatorname{Inf}_i^\mu(p)=\mathbb E_{x\sim\mu}[(p(x)-p(x^{\oplus i}))^2].\]

A junta-free alternative is the \textbf{weak DFKO tail}. Put \(v=\operatorname{Var}_\mu(p)>0\), \(m=\mathbb E_\mu p\), and \(f=(p-m)/\sqrt v\). If \(I=\max_i\operatorname{Inf}_i^\mu(f)\), prove

\[I\le\frac{\eta_0(t)}{t^2}
\quad\Longrightarrow\quad
\Pr_\mu[|f|>t]\ge\eta_1(t)>0,\]

with \(\eta_0,\eta_1\) independent of \(n\), and with

\[\eta_0(t)\ge c\,t^{-A}\]

for fixed positive \(c,A\). This is sufficient, not necessary in general, to obtain AA:

\[\max_i\operatorname{Inf}_i^\mu(p)
\ge \frac14v^2\eta_0(2/\sqrt v)
\ge c'v^{2+A/2}.\]

An exponentially decaying threshold \(\eta_0(t)=e^{-Ct^2}\) is \textbf{not enough} for polynomial AA: it yields only \(v^2e^{-\Theta(1/v)}\).

The simultaneous physical-junta version sought by an earlier route also remains open: find a single low-degree, pointwise-bounded polynomial junta \(q\), on an \(n\)-independent support, achieving both \(\|p-q\|_{L^2(\mu)}^2\le\eta\) and \(\|p-q\|_{L^2(\widetilde{\mu}^R)}^2\le\eta\).

\subsubsection{1B. Established earlier results {[}PROJECT{]}}\label{b.-established-earlier-results-project}

\paragraph{Curie--Weiss and mean-field models}\label{curieweiss-and-mean-field-models}

\begin{itemize}
\tightlist
\item
  \textbf{Curie--Weiss, sharp threshold:} AA holds for \(\beta\le1\), including \(\beta_c=1\), and fails for \(\beta>1\). The route uses exact mixture/slice identities, stability to uniform measure, and a small-\(n\) Poincaré floor.
\item
  \textbf{Verification of critical asymptotics:}
  \(\operatorname{Var}(S)=\Theta(n^{3/2})\) at \(\beta=1\);
  \(\operatorname{Var}(S)/n^{1.5}\to1.164\) and the doubling ratio \(\to2.838\) (predicted \(2^{1.5}=2.828\)).
\item
  The weighted slice distance \(D=\sum_t|p_\beta-p_{\mathrm{Bin}}||t|\) satisfies
  \(D/n^{0.75}\to0.909\), hence \(D/n=O(n^{-1/4})\to0\) and the stability error \(O(d^2D/n)\) vanishes.
\item
  \textbf{Finite-n warning:} the crossover is slow; small-\(n\) numbers do not locate \(\beta_c\). The relevant scale is \(n\gg(\beta-1)^{-3/2}\).
\item
  \textbf{General degenerate mean-field minima:} if the rate function has finite flatness \(I(m)=I(m_*)+c|m|^{2k}\) at its minimum, then \(\operatorname{Var}(S)=\Theta(n^{2-1/k})=o(n^2)\) and the same AA strategy applies. Measured doubling exponents \(1.000,1.500,1.667\) match \(2-1/k=1,1.5,1.667\).
\end{itemize}

\paragraph{Degree-one bounds and tree/mean-field consequences}\label{degree-one-bounds-and-treemean-field-consequences}

For affine \(f\), with finite covariance row bound \(R=\sup_i\sum_j|\operatorname{Cov}(i,j)|\),

\[\max_i\operatorname{Inf}_i(f)\ge \operatorname{Var}(f)^2/R^2.\]

Dobrushin gives \(R\le1/[4(1-\alpha)]\), hence a degree-one bound with \(n\)-uniform constants. Consequences include:

\begin{itemize}
\tightlist
\item
  Dobrushin Gibbs at degree 1 generally;
\item
  one-dimensional Ising at every finite \(\beta\) at degree 1;
\item
  two-dimensional Ising at degree 1 up to \(\beta_c\) (using standard correlation-decay results);
\item
  \(b\)-ary trees at \(\beta_c=\operatorname{artanh}(1/b)\), showing the AA threshold can be \textbf{strictly above} the Dobrushin threshold \(\beta_D\) (and below the KS threshold \(\operatorname{artanh}(1/\sqrt b)\)). Dobrushin is sufficient, not necessary, for AA.
\end{itemize}

\paragraph{Block-independent special cases}\label{block-independent-special-cases}

Correct geometric criterion: flatness requires two support atoms whose \textbf{closed neighbourhoods are disjoint}, equivalently, for two atoms, Hamming distance at least 3. A flip landing off-support still matters because the Dirichlet/flip weight uses \(\lambda(x)>0\).

\begin{itemize}
\tightlist
\item
  Original naive statement ``disconnected support graph'' is false as a criterion. For the \(k=2\) copy law uniform on \(\{00,11\}\), the support graph has no edges, but the two closed neighbourhoods intersect, so no flat nonconstant block function exists.
\item
  For every \(k=2\) law with at least two support atoms, closed neighbourhoods meet, giving positive finite-block Poincaré gap with the unweighted flip form.
\item
  A \(k\ge3\) counterexample remains decisive: \(\lambda\) uniform on \(\{000,111\}\) and \(f=\operatorname{MAJ}_3\) has \(\operatorname{Var}(f)=1/4\) but every single-bit flip influence is zero. Any general block-independent AA claim for unrestricted \(k\ge3\) is therefore false.
\item
  A reported full \(k=2\) theorem has
  \(\max\operatorname{Inf}\ge c\,\delta_{\min}^{-6}\operatorname{Var}^8/d^{12}\),
  with \(\delta_{\min}\) the least block-atom probability. This is accepted in the older log but should be referenced with its exact proof, not casually generalized.
\item
  A stronger \(k=2\) result with law-independent constants was \textbf{accepted only conditionally} because its level-\(d\) comparison lemma was not independently audited. Do not treat it as established.
\end{itemize}

\paragraph{Ising in restricted regimes}\label{ising-in-restricted-regimes}

\begin{itemize}
\tightlist
\item
  \textbf{2D Ising at \(\beta=0.35\), degree 2:} a coarse-graining/SSM proof gives \(n\)-uniform AA, with a formal but practically useless constant around \(10^{-300}\) from block size 100 and huge regularity factors. This reaches beyond \(\beta_D=0.25\) but not \(\beta_c\approx0.4407\).
\item
  \textbf{Higher-temperature expansion:} for sufficiently small \(\kappa=\beta\|J\|_{\rm row}\), connected-cluster counts compare low-degree Walsh moments to product measure. Crude animal counting gives roughly \(\kappa<1/e\) and sharper lattice constants about \(\kappa<1/2.3\). Crucially \(\Delta\) cancels along \((e\Delta)^\ell(\kappa/\Delta)^\ell=(e\kappa)^\ell\). This does not cover the interval up to Dobrushin \(\kappa<1\).
\item
  \textbf{1D Ising, higher degrees:} only a conditional transfer-matrix/LSI route remains in the notes; do not quote it as an established general-degree theorem without the block-routing step.
\end{itemize}

\subsubsection{1C. The junta citation failure and what OZ actually says}\label{c.-the-junta-citation-failure-and-what-oz-actually-says}

An earlier claimed full Dobrushin proof using ``O'Donnell--Zhao 2016 Theorem 3.5 with modified log-Sobolev parameter \(\rho\)'' is \textbf{withdrawn}.

Checked source: O'Donnell--Zhao, \emph{Polynomial bounds for decoupling, with applications}, \href{https://arxiv.org/html/1512.01603v1}{arXiv:1512.01603v1}:

\begin{itemize}
\tightlist
\item
  There is no Theorem 3.5 and no statement of a general-measure MLSI junta theorem with \(\rho\).
\item
  \textbf{Corollary 3.5} concerns uniform Boolean inputs, assumes a small maximum influence of order \(\exp(-Ct^2d^2)\), and uses Gaussian invariance.
\item
  \textbf{Corollary 3.6} is a uniform-measure Fourier-tail theorem: if \(f\) is not an \((\varepsilon,2^{O(k^2/\varepsilon)})\)-junta then a high-level tail has size \(>\exp(-O(k^2)/\varepsilon)\). For bounded degree \(d\), this gives the usual uniform-product junta approximation.
\item
  Ingredients include Walsh orthogonality, commuting independent noise, Gaussian one-block decoupling, and the MOO invariance principle. Glauber MLSI alone does not supply them.
\end{itemize}

\textbf{Do not revive that citation.} Junta approximation over correlated \(\mu\) from a Glauber semigroup MLSI constant remains unproved. ArXiv/DOI links and the checked text are recorded in \texttt{notes/dfko-repair-outcomes.md} and \texttt{.scratch/chk/oz.txt}.

\subsubsection{1D. Important proved structural facts and counterexamples}\label{d.-important-proved-structural-facts-and-counterexamples}

\paragraph{Valid degree-growth example}\label{valid-degree-growth-example}

Heat-bath Glauber on the 3-leaf star \(K_{1,3}\) maps a degree-1 polynomial into degree 3:

\[\mathcal Lx_0=-x_0+c_1(x_1+x_2+x_3)
+c_3x_1x_2x_3,\qquad
c_3=\frac{\tanh(3\beta)-3\tanh\beta}{4}.\]

At \(\beta=0.2\), \(c_3\neq0\). Hence \(P_tx_0\) has a nonzero degree-3 term for small \(t\). Because \(c_3<0\), any small-\(t\) interval must use \(|c_3|\) in its bound.

By contrast, the Ising 3-path preserves affine functions: there \(\mathcal L V_1\subseteq V_1\) and \(P_tV_1\subseteq V_1\). Thus \textbf{do not use the 3-path as a degree-growth counterexample}.

\paragraph{Conditional expectations raise degree: ANOVA is not safe}\label{conditional-expectations-raise-degree-anova-is-not-safe}

The star gives
\[\mathbb E_0[x_0]=\tanh(\beta(x_1+x_2+x_3)),\]
which has a nonzero cubic component. Therefore block-Hoeffding degree, ANOVA degree, and semigroup level are not automatically bounded by original multilinear degree. Any proof claiming block-degree preservation or \(P_t h=e^{-kt}h\) without proof must be rejected.

\paragraph{Bounded versus unbounded MLSI obstruction}\label{bounded-versus-unbounded-mlsi-obstruction}

An unbounded degree-2 example, built from alternating Rademacher pair products, has constant-order variance and requires \(\Omega(n)\) relevant coordinates for a small-error junta approximation. It demonstrates that a pure MLSI assumption cannot, by itself, produce the desired junta theorem. \textbf{It is not a counterexample to bounded Lemma J:} its range grows like \(\sqrt m\), and clipping to \([0,1]\) destroys multilinear degree. The bounded MLSI claim and the bounded Dobrushin claim therefore remain distinct open problems.

\paragraph{Far monomial kills generic gradient localization}\label{far-monomial-kills-generic-gradient-localization}

For uniform product and \(p=x_ix_j\) with \(\operatorname{dist}(i,j)>R\),
\[h_i(x)=(p(x)-p(x^i))^2=x_j.\]

For every \(B(i,R)\)-local \(\ell\),

\[\|h_i-\ell\|_2^2\ge1/4,\]

with optimum \(\ell=1/2\). Moreover \(\sum_kh_k=x_i+x_j\le2\le4d^2\) for \(d=2\), so the GMNR total-influence budget does not rescue localization. Degree-2 monomials can have globally dispersed support. Correlation decay does not imply \(R\)-local approximability.

This is an obstruction to \textbf{local approximation}, not by itself a refutation of TV comparison on the actual support \(\{j\}\).

\paragraph{Correct Dirichlet dynamics for copy blocks}\label{correct-dirichlet-dynamics-for-copy-blocks}

For exact copy \(\lambda\) uniform on \(\{00,11\}\), single-site heat-bath has zero gap: transitions leaving the copy states have rate zero. The unweighted flip form's positive finite-block gap is \textbf{not} the reversible heat-bath MLSI constant. Tensorised block resampling dynamics is the appropriate block-MLSI object. Do not cite a positive single-site Glauber MLSI for the exact copy law.

\paragraph{Bounded global-gradient witness}\label{bounded-global-gradient-witness}

The legal polynomial
\[p_n=x_{i_0}\frac1{n-1}\sum_{j\ne i_0}x_j,\]
is bounded on \(\{0,1\}^n\), has degree 2 and
\[\operatorname{Var}(p_n)=\frac1{16}+\frac1{8(n-1)}.\]
Its \(h_{i_0}\) depends on the other \(n-1\) coordinates. This provides a genuine bounded global example with \(\Theta(n)\) gradient support; no need for the unbounded polynomial \((n^{-1/2}\sum x_j)^2\).

\paragraph{A biased-product mean-square obstruction}\label{a-biased-product-mean-square-obstruction}

A proposed universal comparison \(\|f\|_{L^2(\mu)}\le C\|f\|_{L^2(U)}\) for degree-bounded polynomials is false without a matching first-moment hypothesis. Let \(\mu=\operatorname{Ber}(3/4)^{\otimes n}\) on \(\{\pm1\}^n\), \(U\) uniform, and \(f=1+n^{-1}\sum_i y_i\):
\[\|f\|_{L^2(U)}^2=1+1/n,\qquad
\mathbb E_\mu f^2\to9/4.\]
The mean-square term cannot be omitted.

\paragraph{Duplicate-pair average: no linear variance-to-influence bound}\label{duplicate-pair-average-no-linear-variance-to-influence-bound}

For \(p_m=m^{-1}\sum_{j=1}^m x_{2j-1}x_{2j}\) under uniform independent bits,
\[\operatorname{Var}(p_m)=\frac3{16m},\qquad
\max_i\operatorname{Inf}_i(p_m)=\frac1{2m^2},\]
so \(\max\operatorname{Inf}/\operatorname{Var}=8/(3m)\to0\). This refutes \(n\)-uniform versions of \(\max\operatorname{Inf}\ge c(d)\operatorname{Var}\); a polynomial in \(v\), here \(v^2\), is necessary.

\subsubsection{1E. Accepted partial analytic steps (not AA proofs)}\label{e.-accepted-partial-analytic-steps-not-aa-proofs}

\paragraph{\texorpdfstring{Corrected \(L^2\) influence ledger}{Corrected L\^{}2 influence ledger}}\label{corrected-l2-influence-ledger}

For bounded \(p,q\), with \(C_{\rm RN}\le(1-\delta_0)/\delta_0\),

\[\mathbb E_\mu[(\Delta_i(p-q))^2]
\le2(1+C_{\rm RN})\|p-q\|_{2,\mu}^2.\]

This follows by flipping variables and using \(\mu(x^i)/\mu(x)\le C_{\rm RN}\). It gives a per-coordinate influence comparison. Separately, \(|\operatorname{Var}_\mu(p)-\operatorname{Var}_\mu(q)|\le4\|p-q\|_2\).

\textbf{Requirements:} preserve the distinction between \(\eta\) and \(\sqrt\eta\), and require \(q\) to be both multilinear degree \(\le d\) and pointwise bounded. Clipping a general junta to \([0,1]\) need not preserve multilinear degree and can invalidate the product-AA application.

GMNR Lemma 2.11 correctly gives \(\sum_i(f(x)-f(x^i))^2\le4d^2\) pointwise for bounded degree-\(d\) \(f\), under any distribution. Earlier denials of that citation were themselves wrong and have been withdrawn.

\paragraph{Joint endpoint functional TV lemma}\label{joint-endpoint-functional-tv-lemma}

For ordered pairs \((x^R,y^R)=(F^R(z),F^R(z^k))\) and \((x,y)=(F(z),F(z^k))\), on \(A_S=G_S\cap G'_S\), both truncations agree with their untruncated counterparts. For bounded \(\Phi\),

\[\left|\mathbb E_{\pi_k}\Phi(x_S^R,y_S^R)-\mathbb E_{\pi_k}\Phi(x_S,y_S)\right|
\le 2C_0M_Se^{-c_0\delta R}+2L2^{-L}.\]

The path functional is correctly a function of \textbf{both} endpoints; the previous definition from only \(x^R\) was invalid.

\textbf{Still missing:} the additional domination/aggregation connecting that functional to \(\operatorname{Inf}_i^\mu\). This is not an influence-transfer theorem or AA proof.

\subsubsection{1F. Conditional cumulant/MGF reduction {[}CONDITIONAL{]}}\label{f.-conditional-cumulantmgf-reduction-conditional}

Suppose, for suitably normalized \(f\), both a hypercontractive moment and cumulant estimate hold:

\[\|f\|_q\le(K_0q)^{K_1d}\|f\|_2, \tag{HC}\]

\[|\kappa_q(f)|\le(C_1dq)^{C_2dq}I^{q/2-1}q!. \tag{CU}\]

Then over the relevant interval,
\[\log M(s)=s^2/2+O(s^3\sqrt I).\]

At \(s=t+C\) for a suitable constant \(C\), the tilted mean lies above \(t\) with bounded tilted variance, so Chebyshev gives constant central-mass probability. Backward tilting gives

\[\Pr[f>t]\ge c\exp(-t^2/2+O(t^3\sqrt I)).\]

Using both signs and arranging \(t^3\sqrt I=O(1)\) gives a weak DFKO tail with a polynomial-in-\(t\) influence threshold, and hence polynomial AA conditionally via the usual cutoff.

\textbf{HC+CU remain completely unproved.} Dirichlet proofs, derivative/ANOVA versions, and the small-\(\kappa\) corollary are not accepted. This is a useful sufficient package only, not a completed route.

\subsubsection{1G. Exact hybrid/Taylor framework and remaining covariance gap}\label{g.-exact-hybridtaylor-framework-and-remaining-covariance-gap}

With \(Z^{(i)}=(Y_1,\dots,Y_i,X_{i+1},\dots,X_n)\) under \(\mu\otimes\nu\), the telescoping and third-order Taylor identities are exact. They define a framework for dependent replacement, not a proof of invariance. The needed first-order bound is

\[\left|\sum_i\mathbb E\!\left[\psi'(f^{(i)})(X_i-Y_i)\partial_i p/\sqrt v\right]\right|
\le C\|\psi'\|_\infty I_\mu^c,\]

with \(C,c\) independent of \(n\). This remains open. Even if proved, the smooth \(C^3\) error \(\varepsilon^{-3}I^{1/2}\) is only polynomial in \(t^{-1}\) and does not beat the exponential DFKO tail for large \(t\). A separate non-smooth/exponential tail comparison is required.

Upper moment bounds yield only \textbf{upper} tails, never a lower large-\(t\) bound \(\Pr[|f|>t]\ge ct^{-B}\).

\subsubsection{1H. Q4 / fourth-moment route: status and warning}\label{h.-q4-fourth-moment-route-status-and-warning}

A proposed minimal lemma Q4(\(d\)) asks whether, for degree-\(\le d\) \(p\) under Dobrushin,

\[\|p-\mathbb E_\mu p\|_{L^4(\mu)}
\le C_d\|p-\mathbb E_\mu p\|_{L^2(\mu)},\]
with \(C_d\) independent of \(n\) and \(p\).

\textbf{It is unproved, as are the derivative/ANOVA variants and the general induction.} The proposed Riesz estimate \(\|p-P_{t_*}p\|_4\le\frac12\|p\|_4\) was only a sketch and is not supplied. Conditional expectations can increase degree, so an induction cannot freely assume preservation.

If Q4 held at all degrees and with all derivative variants, then OZ's algebraic sup-norm bound plus Dobrushin Poincaré and reverse Hölder would give the displayed AA bound. But this does not constitute a simplification: Q4 encapsulates a very strong hypercontractive/DFKO-type statement and may be as difficult as the main target. Bonami and nonhomogeneous OZ constants carry exponential \(d\) factors, so calling the \(v^2\) result ``sharper'' without tracking them would be misleading.

\subsubsection{1I. Principal open lemmas due to current gaps}\label{i.-principal-open-lemmas-due-to-current-gaps}

\begin{itemize}
\tightlist
\item
  Full Dobrushin AA, especially \(d\ge2\) beyond established small-\(\kappa\) regions. In particular, general degrees in the 2D Ising window \((\beta_D\approx0.25,\beta_c\approx0.4407)\) and the Dobrushin interval after the animal threshold remain open.
\item
  The polynomial-threshold weak DFKO tail described above.
\item
  The simultaneous low-degree bounded junta \(q\) under both \(\mu\) and \(\widetilde{\mu}^R\) (``C1'').
\item
  The intermediate path-TV comparison and subsequent influence domination (``C2'').
\item
  HC+CU of §1F, including the Dirichlet-form control.
\item
  The first-order hybrid covariance estimate of §1G.
\item
  Correct block/ANOVA degree control, and the variance and argmax expressions:
  \[v=\mathbb E_\tau v_\tau+\operatorname{Var}_\tau m_\tau,\qquad
  \max_i\mathbb E_\tau\phi_i\le\mathbb E_\tau\max_i\phi_i,\]
  with no assertion of equality in the latter.
\item
  Any smaller-scale replacement with a genuinely independent mechanism.
\end{itemize}

\begin{center}\rule{0.5\linewidth}{0.5pt}\end{center}
